\chapter{Giới thiệu}\label{chap:gioithieu}
\section{Đặt vấn đề}
Wikipedia thay đổi liên tục khi người dùng sửa bài, tạo trang và thực hiện các thao tác khác. Khi lượng sự kiện lớn, người xem khó nhận ra trang nào đang thu hút hoạt động và chỉnh sửa nào nên được kiểm tra trước. Luồng thay đổi công khai của Wikimedia cung cấp dữ liệu để trả lời hai câu hỏi đó.

Xử lý theo thời gian thực gặp các khó khăn: sự kiện có thể đến muộn, lặp lại khi kết nối gián đoạn hoặc chứa dữ liệu không hợp lệ. Hệ thống cần tiếp tục phục vụ kết quả khi có thay đổi, đồng thời bảo vệ thông tin cá nhân và giải thích rõ ý nghĩa của số liệu.
\section{Mục tiêu và phạm vi}
WikiPulse tiếp nhận dữ liệu từ Wikimedia hoặc phát lại dữ liệu mẫu, tổng hợp hoạt động theo thời gian, xếp hạng xu hướng và phát hiện các dấu hiệu cần ưu tiên xem xét. Dashboard cung cấp thông tin cho người tuần tra và người quan tâm tới Wikipedia.

Phiên bản được nghiên cứu chạy trên một node, hỗ trợ hai phiên bản Wikipedia tiếng Anh và tiếng Việt theo cấu hình, lưu dữ liệu cục bộ hoặc trên HDFS. Việc xác thực, phân quyền, high availability, mô hình học máy và triển khai cloud cho production nằm ngoài phạm vi; nhóm có thực hiện một phép thử tạm trên cloud. Điểm ưu tiên không phải kết luận về hành vi phá hoại.
\section{Tổ chức thực hiện}
\begin{table}[H]
\centering\small
\caption{Phân công công việc}\label{tab:phancong}
\begin{tabularx}{\linewidth}{@{}lX@{}}\toprule
\textbf{Thành viên} & \textbf{Phụ trách} \\ \midrule
Trần Quang Trọng & Tích hợp hệ thống, môi trường triển khai và kiểm thử chung \\
Lê Nguyễn Phước Thành & Thu thập sự kiện, phát lại và kết nối nguồn dữ liệu \\
Nguyễn Ngọc Linh & Xử lý luồng, chuẩn hóa, tổng hợp và tính điểm \\
Hoàng Văn Nhân & Lưu trữ phục vụ truy vấn và API \\
Nguyễn Đức Mạnh & Dashboard và kiểm thử chấp nhận \\
\bottomrule\end{tabularx}\end{table}
\section{Cấu trúc báo cáo}
Các chương tiếp theo trình bày nền tảng công nghệ, các hướng tiếp cận liên quan, thiết kế hệ thống, kết quả kiểm thử và những hướng phát triển tiếp theo.
